\lecture{12}{Sep 17 10:15}{How Many Solutions Are There?}

\section{System Solution Analysis}

\subsection{Consistency}
\begin{theorem}[Consistency Condition]\label{thm:consistency-condition}
    A linear system is consistent if and only if its augmented coefficient matrix is row equivalent to a matrix in reduced row echelon form that does not contain a type-B row.
\end{theorem}

\begin{proof}
    Form the augmented coefficient matrix of the system and run the Gauss-Jordan elimination algorithm.
    \begin{itemize}
        \item All operations during the algorithm are elementary row operations
        \item All matrices (initial, final, and intermediate) are row equivalent to the original matrix
        \item All linear systems represented by these matrices are equivalent
        \item If at any point a type-B row appears, the system is inconsistent
        \item If the algorithm completes, the output is a RREF matrix without a type-B row
        \item This matrix represents a system equivalent to the original system
        \item A solution is available by back substitution
    \end{itemize}
\end{proof}

\subsection{Number of Solutions}
Let $A$ be a linear system with $n$ unknowns over a field $F$, and suppose the augmented coefficient matrix is row equivalent to a RREF that does not contain a type-B row.

\begin{enumerate}
    \item If the RREF matrix has exactly $n$ non-zero rows, then the system has a \textbf{unique solution}
    \item If the RREF matrix has $k < n$ non-zero rows, then there are $n - k$ free variables:
    \begin{itemize}
        \item If field $F$ has infinitely many scalars, there are \textbf{infinitely many solutions}
        \item If field $F$ has $N$ different scalars, there are $N^{n-k}$ solutions
    \end{itemize}
\end{enumerate}

\begin{note}
    The RREF matrix has at most $n$ non-zero rows because each non-zero row has a pivot in one of the first $n$ columns, and everything else in a pivot column is zero.
\end{note}

\subsection{Homogeneous Systems}

\begin{note}
    A homogeneous system is always consistent because the trivial solution always exists.
\end{note}

\begin{theorem}[Non-trivial Solutions]\label{thm:nontrivial-solutions}
    A homogeneous system with more unknowns than equations has a non-trivial solution.
\end{theorem}

\begin{proof}
    Let $A$ be a homogeneous system with $m$ equations in $n$ unknowns where $m < n$.
    \begin{itemize}
        \item The RREF of the coefficient matrix has no type-B rows (system is homogeneous)
        \item There are at most $m$ pivots in the RREF matrix (at most one per row)
        \item Therefore, there are at most $m$ restricted variables
        \item Therefore, there is at least one free variable since $n - m > 0$
        \item Any choice of $t \in F$ for the free variable leads to a solution that is an $n$-tuple with $t$ as one component
        \item Including the choice $t = 1$
        \item Since $1 \neq 0$ by axiom viii in Definition \ref{def:field}, there exists a non-trivial solution
    \end{itemize}
\end{proof}